\section{A second result: a majorant for the complete Pell radical}
\label{sec:full-radical}

A separate construction improves the lowest inspected degree with nine
strictly positive unknowns. Its degree is
\[
 \boxed{1\,158\,228},
\]
compared with package~71's $1\,379\,730$, a reduction of $16.0540\%$.
Its delivered evaluator uses $726$ binary operations, compared with $478$;
this is a degree--size tradeoff. All four coefficient payloads are the
coherent original31 tuple. The new polynomial does not add a witness or
change the native compiler.

\subsection{The complete radical and the denominator that can be cleared}

For $A\ge3$, $C,i\ge1$, and $t\ge3$, put
\begin{align*}
 D&=(A^2-1)C^2+1,& e&=i(4A-5)C,\\
 F&=(D-1)e^2D^2+1,& L&=CD+1-A,\\
 G&=A+FL,& H&=C+Ft,\\
 I&=(G^2-1)H^2+1.
\end{align*}
All these registers are positive, and $G>1$. Regard $A$ as a formal
indeterminate over $\mathbb Q(C,i,t)$. The radical $\sqrt{DFI}$ has a unique
formal Laurent expansion with leading term
$2^{10}C^{24}i^5tA^{23}$. Its polynomial part provides the construction
\begin{equation}
 T_0=2^{16}C^7[A^{\ge0}]\sqrt{DFI}.
 \label{eq:full-root-definition}
\end{equation}
The expression $T_0$ is an integer polynomial with $654$ monomials. Its
complete table is the \texttt{majorant} field of
\path{nine_full_radical/laurent_candidate.json}; the coordinate order is
$(A,C,i,t)$. This table is an ordinary division-free definition.
Formula~\eqref{eq:full-root-definition} explains how it was discovered.

The exact polynomial
\begin{equation}
 \Delta_0=T_0^2-2^{32}C^{14}DFI
 \label{eq:full-root-defect}
\end{equation}
has $A$-degree $22$. Together with the leading term of $T_0$, this also
proves the formal-polynomial-part assertion: dividing the difference of
the squares by their formal sum leaves only negative powers of $A$.
The factor $C^7$ clears variable denominators without raising the
$A$-degree.

\paragraph{Complete positivity certificate.}
Under
\[
 A=x+3,\qquad C=y+1,\qquad i=z+1,\qquad t=u+3,
\]
$T_0$ has exactly $9\,216$ nonzero coefficients and $\Delta_0$ exactly
$32\,637$. Every one is a positive integer, including the constant
coefficients. The archive contains both complete expansions and an
independent verifier that reconstructs the radical and checks every
coefficient. Hence, throughout the stated real domain,
\begin{equation}
 T_0>2^{16}C^7\sqrt{DFI}>0.
 \label{eq:full-root-bound}
\end{equation}
This is a finite algebraic proof. No tests on a bounded selection of
integer inputs replace the coefficient identities.

The improvement comes from approximating the whole radical. The preceding
seventh-order construction approximated $\sqrt{DF}\,G$ and multiplied the
error by the high-degree factor $H$. Taking the polynomial part of the
complete radical prevents that multiplication of the approximation error.

\subsection{The final polynomial and its correctness}

Retain the source and Pell registers of package~71, with positive unknowns
\[
 b,g,h,s,w,C,i,j,n.
\]
The source has convolution height $1873$ and maximum coordinate weight
$116$. Its full literal circuit, all kernels, and the fixed coefficient
generator and compressed integer payloads are included in the companion
directory. In particular,
\begin{gather*}
 U=N^2w,\quad Y=N^2s,\quad M=RY,\quad P=2M^2U,\\
 K=R+1+h(P-1),\quad A=M(U+1),\quad B_1=2R+1,\quad W_1=bw,\\
 J=4A-5,\quad Q=3W_1C-2(W_1^2-1),\quad t=B_1+(j-1)C,\\
 a_0=(P^2-1)K^2+1,\quad V_0=PK,\\
 m=K^2(4\operatorname{cap}-g)
      -16\operatorname{cap}(C-KY)^2,\\
 \Omega=Jm,\qquad E_0=J(n+W_1^2)+Q.
\end{gather*}
Define
\begin{gather*}
 S=2^{18}C^7,\qquad T=4T_0(A,C,i,t),\qquad a_1=S^2DFI,\\
 \Theta=E_0-\Omega(V_0+T),\\
 \boxed{\displaystyle
 P_9=(\Theta^2-\Omega^2(a_0+a_1))^2
                -4\Omega^4a_0a_1.}
\end{gather*}
All displayed registers are abbreviations, not additional unknowns.
The authoritative evaluator is
\path{nine_full_radical/laurent9_positive.json}.

The unchanged source estimates give $R\ge1$, $A\ge3$, and hence
$t\ge3$, unconditionally on positive inputs. Thus
$T>\sqrt{a_1}$ before any source zero has been decoded. Also
$V_0^2-a_0=K^2-1>0$. If $m=0$ the norm is $E_0^4>0$; otherwise its real
factorization gives signed roots $u_0^2=a_0$, $u_1^2=a_1$ and
\[
 E_0/J=m(V_0+T+u_0+u_1).
\]
The parenthesis is strictly positive. Therefore $m>0$ before the
square-root integrality argument. The inherited margin argument gives
$C>K$, $C>B_1+W_1$, and $DFI>a_0$. In particular $a_1>a_0$ and the
rational sum $u_0+u_1$ is nonzero. The identity
\[
 2(u_0+u_1)u_0=(u_0+u_1)^2+a_0-a_1
\]
makes both roots rational and then integral. Since $(u_1/S)^2=DFI$,
the rational-root argument also makes $u_1/S$ an integer. The factor
equation recovers $J\mid Q$, and the existing ratio theorem supplies the
encoded positive source zero.

Conversely, choose the inherited ratio witnesses from a source zero,
take the positive integral roots $r_0=\sqrt{a_0}$ and
$r_1=\sqrt{DFI}$, and set
\[
 n=m(V_0+T+r_0+Sr_1)-W_1^2-Q/J.
\]
It is an integer. The inherited bound
$W_1^2+Q/J<2D$, together with $T>S\sqrt{DFI}>SD>2D$ and $m\ge1$,
shows $n>0$. One norm factor now vanishes. Thus the positive existence
statements are equivalent in both directions. Combining with the
specified native compiler and refutation calibration gives the desired
ZFC-provable consistency equivalence. These inherited semantic results
are not being presented as newly formalized proofs.

\subsection{Exact degree and evaluator cost}

For this unchanged source let
\[
 r=\deg R=7963,\quad q=\deg N=2115,
 \quad a=\deg A=16425,\quad \omega=\deg\Omega=82360.
\]
The leading $T_0$ monomial is $2^{26}A^{23}C^{31}i^5t$; the highest
weighted monomial of~\eqref{eq:full-root-defect} is
\[
 141664530202624\,A^{22}C^{42}i^{10}t^2.
\]
Therefore
\[
 \deg T=23a+r+36=385774,\qquad
 \deg(T^2-a_1)=22a+2r+52=377328.
\]
Write $\delta_0=K^2-1$, $\delta_1=T^2-a_1$, and
\[
 B_* =E_0^2-2E_0\Omega(V_0+T)
                        +\Omega^2(\delta_0+\delta_1).
\]
The exact identity
\[
 P_9=B_*^2+4\Omega^2V_0TB_*
                  +4\Omega^4(\delta_0T^2+a_0\delta_1)
\]
has four candidate degrees
\[
 1084096,\qquad1149781,\qquad1158228,\qquad821246.
\]
The unique dominant term is $4\Omega^4\delta_0T^2$.
On the positive diagonal its leading coefficient is
\[
 2^{166}a_b^{578828}c_b^{276}>0,
 \qquad a_b=2^{2134939},\quad c_b=2^{17079501}.
\]
Consequently the asserted degree is attained by the actual fixed
coefficient instance.

The $726$-operation evaluator uses a $110$-term correction to the old
majorant, in the existing registers $J,C,i,t,D$:
\[
 T=64C^7(T_7-640i-4096)H
                   -\operatorname{Corr}(J,C,i,t,D).
\]
The correction table and the exact identity check are included. There
are $430$ multiplications, $265$ additions, and $31$ subtractions;
coefficient leaves and register reuse are free under the same convention
as the comparison source.

The new degree limit is now set by the first-root defect $K^2-1$. A tighter
simultaneous bound for both norm roots is a concrete further research
question, as are cheaper evaluators of the correction and an application
to the ten-integer finalizer. None of those additional improvements is
claimed here.
